% !TEX program = pdflatex
% THEME 2 — Difficile — Corrigé
\documentclass[11pt,a4paper]{article}
\usepackage[utf8]{inputenc}
\usepackage[T1]{fontenc}
\usepackage{lmodern}
\usepackage[french]{babel}
\usepackage{amsmath,amssymb,mathtools}
\usepackage{icomma}
\usepackage{siunitx}
\sisetup{output-decimal-marker={,},per-mode=symbol}
\DeclareSIUnit{\atm}{atm}
\DeclareSIUnit{\bar}{bar}
\DeclareSIUnit{\mmHg}{mmHg}
\usepackage[version=4]{mhchem}
\usepackage{xcolor}
\usepackage{colortbl}
\usepackage[most]{tcolorbox}
\usepackage{enumitem}
\usepackage{array}
\usepackage{booktabs}
\usepackage{fancyhdr}
\usepackage[a4paper,margin=2.1cm,headheight=26pt,headsep=10pt]{geometry}

\definecolor{primary}{RGB}{150,98,162}
\definecolor{primarydark}{RGB}{92,52,110}
\definecolor{excol}{RGB}{0,135,90}
\definecolor{attncol}{RGB}{200,40,55}

\newtcolorbox{rep}[1][]{enhanced,breakable,colback=excol!5,colframe=excol,
  boxrule=0.8pt,arc=2pt,left=3mm,right=3mm,top=2mm,bottom=2mm,
  fonttitle=\bfseries,coltitle=white,title={#1}}

\pagestyle{fancy}\fancyhf{}
\fancyhead[L]{\small\textcolor{primarydark}{\textbf{Fiche 7} — Loi des gaz}}
\fancyhead[R]{\small\textcolor{primarydark}{\textsc{Thème 2 · Difficile · Corrigé}}}
\fancyfoot[C]{\small\thepage}
\renewcommand{\headrulewidth}{0.8pt}
\renewcommand{\headrule}{\hbox to\headwidth{\color{primary}\leaders\hrule height \headrulewidth\hfill}}
\setlength{\parindent}{0pt}\setlength{\parskip}{3pt}

\newcounter{exo}
\newenvironment{cor}[1][]{\refstepcounter{exo}\medskip
  \noindent{\color{primarydark}\bfseries Corrigé \theexo.}\ \textit{#1}\par\nopagebreak\smallskip}{\par}
\newcommand{\rf}[1]{\textbf{\textcolor{attncol}{#1}}}

\begin{document}

\begin{center}
{\Large\bfseries\color{primarydark} Thème 2 — Niveau Difficile — Corrigé}
\end{center}
\textcolor{primary}{\hrule height 1pt}
\medskip

\begin{cor}[Bouteille d'azote comprimé]
\textbf{(a)} $T_1=25+273=\qty{298}{\kelvin}$, puis $n=\dfrac{p_1V}{RT_1}$ :
\[ n=\frac{\qty{10}{\atm}\times\qty{20}{\litre}}{\num{0.082}\times\qty{298}{\kelvin}}
   =\frac{200}{24{,}44}=\rf{\qty{8.2}{\mole}}. \]
\textbf{(b)} $m=n\,M(\ce{N2})=8{,}2\times28\approx\rf{\qty{229}{\gram}}$.

\textbf{(c)} $T_2=100+273=\qty{373}{\kelvin}$.\\
\emph{(i)} avec $pV=nRT$ ($n$ et $V$ inchangés) :
$p_2=\dfrac{nRT_2}{V}=\dfrac{8{,}2\times0{,}082\times373}{20}=\dfrac{250{,}8}{20}=\rf{\qty{12.5}{\atm}}$.\\
\emph{(ii)} avec la loi isochore :
$p_2=p_1\dfrac{T_2}{T_1}=10\times\dfrac{373}{298}=\rf{\qty{12.5}{\atm}}$.\\
Les deux méthodes donnent \rf{le même résultat} : la loi isochore n'est que $pV=nRT$ avec $V$ et $n$ figés.

\textbf{(d)} En chauffant, la pression interne augmente (ici de $10$ à $\qty{12.5}{\atm}$) : au-delà d'une
certaine valeur, la bouteille risque de \rf{céder ou d'exploser}. On stocke donc les gaz comprimés au frais.
\end{cor}

\begin{cor}[Pneu qui fuit]
\textbf{(a)} $n_1=\dfrac{p_1V}{RT}=\dfrac{2{,}5\times3{,}0}{0{,}082\times300}=\dfrac{7{,}5}{24{,}6}
=\rf{\qty{0.305}{\mole}}$ ; \quad
$n_2=\dfrac{p_2V}{RT}=\dfrac{2{,}0\times3{,}0}{24{,}6}=\dfrac{6{,}0}{24{,}6}=\rf{\qty{0.244}{\mole}}$.\\
\textbf{(b)} Air échappé : $\Delta n=n_1-n_2=0{,}305-0{,}244=\rf{\qty{0.061}{\mole}}$.\\
\textbf{(c)} Fraction perdue : $\dfrac{\Delta n}{n_1}=\dfrac{0{,}061}{0{,}305}=0{,}20=\rf{\qty{20}{\percent}}$.
On retrouve directement ce résultat par $1-\dfrac{p_2}{p_1}=1-\dfrac{2{,}0}{2{,}5}=0{,}20$ : à $V$ et $T$
constants, la quantité de matière est proportionnelle à la pression.\\
\textbf{(d)} La variable qui change est \rf{$n$} (la quantité de gaz), et non $p$, $V$ ou $T$
« librement ». La baisse de pression n'est que la \emph{conséquence} de la perte de matière : ce n'est
pas une transformation à $n$ constant, contrairement à celles du Thème 1.
\end{cor}

\begin{cor}[Identifier un gaz inconnu]
\textbf{(a)} $T=27+273=\qty{300}{\kelvin}$, puis $n=\dfrac{pV}{RT}$ :
\[ n=\frac{\qty{1.0}{\atm}\times\qty{2.46}{\litre}}{\num{0.082}\times\qty{300}{\kelvin}}
   =\frac{2{,}46}{24{,}6}=\rf{\qty{0.10}{\mole}}. \]
\textbf{(b)} $M=\dfrac{m}{n}=\dfrac{\qty{4.4}{\gram}}{\qty{0.10}{\mole}}=\rf{\qty{44}{\gram\per\mole}}$.

\textbf{(c)} Une masse molaire de $\qty{44}{\gram\per\mole}$ oriente vers le \rf{dioxyde de carbone \ce{CO2}}
($12+2\times16=44$) ou vers le \rf{propane \ce{C3H8}} ($3\times12+8=44$). \rf{Non}, la masse molaire seule
ne suffit pas à trancher : deux gaz différents peuvent avoir la même masse molaire. Il faudrait une
information supplémentaire (test chimique, spectre\dots) pour conclure avec certitude.
\end{cor}

\end{document}
